\documentclass[11pt]{article}
\usepackage[a4paper,margin=28mm]{geometry}
\usepackage[T1]{fontenc}
\usepackage{lmodern,microtype}
\usepackage{amsmath,amssymb,amsthm,mathtools}
\usepackage{booktabs,array,longtable,tabularx}
\usepackage{placeins}
\usepackage{enumitem}
\usepackage[dvipsnames]{xcolor}
\usepackage{hyperref}
\hypersetup{colorlinks=true,linkcolor=MidnightBlue,citecolor=MidnightBlue,urlcolor=MidnightBlue,
 pdftitle={Fast distributed edge coloring by local palette selection},pdfauthor={}}
\usepackage[nameinlink,noabbrev]{cleveref}
\setlist{itemsep=2pt,topsep=5pt}
\setlength{\emergencystretch}{2em}
\newtheorem{theorem}{Theorem}[section]
\newtheorem{lemma}[theorem]{Lemma}
\newtheorem{corollary}[theorem]{Corollary}
\newtheorem{proposition}[theorem]{Proposition}
\theoremstyle{definition}
\newtheorem{definition}[theorem]{Definition}
\theoremstyle{remark}
\newtheorem{remark}[theorem]{Remark}
\newcommand{\N}{\mathbb N}
\newcommand{\F}{\mathcal F}
\newcommand{\eps}{\varepsilon}
\newcommand{\poly}{\operatorname{poly}}
\newcommand{\ceil}[1]{\left\lceil #1\right\rceil}
\newcommand{\floor}[1]{\left\lfloor #1\right\rfloor}
\newcommand{\abs}[1]{\left|#1\right|}
\newcommand{\set}[1]{\left\{#1\right\}}
\newcommand{\logstar}{\log^*}
\newcommand{\LOCAL}{\ensuremath{\mathsf{LOCAL}}}
\newcommand{\CONGEST}{\ensuremath{\mathsf{CONGEST}}}
\DeclareMathOperator{\degmax}{\Delta}
\title{Fast distributed edge coloring\\by local palette selection}
\author{}
\date{30 September 2026}

\begin{document}
\maketitle
\begin{center}
\fbox{\begin{minipage}{0.91\linewidth}
\textbf{Highlight.} A graph of maximum degree at most a known bound
\(\Delta\), supplied with a proper vertex $2$-coloring, admits a
\emph{deterministic proper edge coloring with $O(\Delta)$ colors in
$O(1+\logstar\Delta)$ communication rounds}. No vertex identifiers are
needed, and the bound is independent of the number of vertices.
\end{minipage}}
\end{center}

\begin{abstract}
We prove the highlight theorem by an elementary construction of expanding
sets and two applications of Hall's marriage theorem. Each old edge color
specifies a short list of new colors. One side of the bipartition proposes
disjoint sublists, and the other side accepts distinct representatives.
This reduces a palette of $L\Delta$ colors to $O(\Delta\log L)$ colors.
We also prove a sharper $O(\Delta\log L/\log\log L)$ reduction, an explicit
$\lfloor16.5\Delta\rfloor$ final palette, and a direct two-round
$O(\Delta^{5/3})$ algorithm. For arbitrary edge lists, a related construction
colors lists of size $O(\Delta\log\Delta/\log\log\Delta)$ in three rounds;
batching and partial coloring give linear-list and degree-plus-one-list
tradeoffs. We give general-graph consequences, a tight one-round lower
bound under explicit initialization conventions, and limitations of two
particular set-system constructions. All new constructions and reductions
are proved, including their numerical estimates. The probabilistic method
is used to construct deterministic finite tables, and no efficiency of
their local construction is assumed.
\end{abstract}

\clearpage
{\small\tableofcontents}
\clearpage

\section{Results, model, and perspective}\label{sec:overview}

The input to the central result is a bipartite graph \emph{with its two
sides supplied}. We call the sides black and white. This information is
stronger than a promise of bipartiteness. The result applies to irregular
graphs as well as regular graphs; a common upper bound $\Delta$ on their
degrees is assumed known. The empty graph is immediate, so throughout the
proofs we may assume $\Delta\ge1$.

\subsection{Communication and colors}
In the synchronous \LOCAL\ model, each vertex can send an arbitrary finite
message along each incident edge in each round and perform arbitrary finite
local computation between rounds. Each vertex initially knows its degree,
its own distinct ports $0,\ldots,\deg(v)-1$, and its side of the supplied
bipartition. Opposite-end ports are initially unknown. On general graphs
we use distinct identifiers from a universe of polynomial size in $n$,
the number of vertices. For these general-graph upper bounds, $n$ (or a polynomial estimate of it and of the identifier range) may be supplied. The anonymous \LOCAL\ bounds permit any common upper bound $\Delta$. For general-graph and bandwidth statements we assume $1\le\Delta\le n-1$ on a nonempty simple graph; this holds when $\Delta$ is the actual maximum degree. In the general-graph setting a larger supplied bound may be clamped using the supplied $n$. Unless an intermediate one-sided state is explicitly
specified, \emph{both endpoints output the same final color of their edge}.
These conventions are especially relevant to exact one- and two-round bounds.

An edge coloring is proper if incident edges receive different colors.
In list edge coloring, each edge $e$ has a finite set $L_e$ of permissible
colors, initially known at both endpoints. Labels have finite encodings
and a common total order; their lengths and their global universe are
unrestricted. All list results and the ordinary corollaries derived from them are in \LOCAL. The ordinary palette-reduction algorithms and the direct two-round construction also fit
\CONGEST, where each edge carries $O(\log n)$ bits per round, under the
same unrestricted-local-computation convention; see
\cref{sec:implementation}. Graphs are finite and simple unless otherwise
stated. For a simple graph $G$, write
\[
 d_G(e)=\deg_G(u)+\deg_G(v)-2\qquad(e=uv)
\]
for the degree of $e$ in the line graph. Thus a degree-plus-one edge list
has size at least $d_G(e)+1$. We use ``degree-plus-one'' in this line-graph
sense, not to mean $\Delta+1$.

All tables below are uniform sufficient guarantees, not optimality claims.
Logarithms in asymptotic expressions are cut off at an absolute constant.
For exact parameter choices, put
\begin{equation}\label{eq:rho}
 h=1+\logstar\Delta,\qquad \ell=1+\logstar n,\qquad
 \rho=\min\{r\in\N:r\ge3,\ r^{r-2}\ge720\Delta^4\}.
\end{equation}
Here $\logstar$ uses base two, and
$\rho=\Theta(\log\Delta/\log\log\Delta)$ as $\Delta\to\infty$.
Taking logarithms of its defining inequality proves this last assertion:
$(r-2)\ln r$ must be $\Theta(\ln\Delta)$, and constant multiples of
$\ln\Delta/\ln\ln\Delta$ bracket the least feasible $r$ for large $\Delta$.

\subsection{Palette size versus time}
\begin{table}[htbp]
\centering\small
\begin{tabularx}{\linewidth}{@{}lXl@{}}
\toprule
Task with a supplied bipartition & Sufficient palette or list size & Rounds\\
\midrule
Ordinary & $\Delta^2$ & $1$\\
Ordinary & $<3\Delta^{5/3}$ & $2$\\
Ordinary & $O\!\left(\Delta\frac{\log^{(t)}\Delta}{\log^{(t+1)}\Delta}\right)$,
 any fixed $t\ge1$ & $t+2$\\
Ordinary & $\lfloor16.5\Delta\rfloor$ & $O(h)$\\
Lists & $\Delta^2$ & $2$\\
Lists & $\rho\bigl(d_G(e)+1\bigr)$ on edge $e$ & $3$\\
Lists & $2\Delta-2+\rho(2k-1)$, $1\le k\le\Delta$ & $O(h+\Delta/k)$\\
Lists & $\ceil{(2+\eps)\Delta}$, $\eps>0$ & $O(h+1+\rho/\eps)$\\
Lists & $d_G(e)+1$ on edge $e$ & $O(\rho(1+\log\Delta))$\\
\bottomrule
\end{tabularx}
\caption{Deterministic bounds with a supplied bipartition. Iterated logarithms
are denoted by $\log^{(t)}$. The $t+2$ bound pipelines reductions by alternating
the active side, with a final notification to both endpoints.}
\label{tab:bipartite}
\end{table}

In particular, $100\Delta$-element lists take $O(\rho)$ rounds, whereas
ordinary $O(\Delta)$ coloring takes $O(h)$ rounds. Choosing
$k=\ceil{\Delta/h}$ in the batching bound gives lists of size
$O(\Delta\rho/h)$ in $O(h)$ rounds. Choosing
$k=\ceil{\Delta/(1+\log\Delta)}$ gives
\begin{equation}\label{eq:neartwo}
 2\Delta+O\!\left(\frac{\Delta}{\log\log\Delta}\right)
 \quad\text{colors or list entries in }O(1+\log\Delta)\text{ rounds}.
\end{equation}
Exact $2\Delta-1$ colors follow by assigning every edge the same list of
that size. Its time is $O(\rho(1+\log\Delta))$.

\begin{table}[htbp]
\centering\small
\begin{tabularx}{\linewidth}{@{}lXl@{}}
\toprule
Task on a general graph & Sufficient size & Total rounds\\
\midrule
Ordinary & $3\lfloor16.5\Delta\rfloor\le50\Delta$ & $O(\ell)$\\
Lists & $O_c(\Delta^2\log\Delta/2^{c h})$, any fixed $c>0$ & $O_c(\ell)$\\
Lists & $100\Delta$ & $O(1+\log\Delta+\ell)$\\
Ordinary or lists & $2\Delta+O(\Delta/\log\log\Delta)$ & $O(1+\log\Delta+\ell)$\\
Lists & $d_G(e)+1$ on edge $e$ & $O(\rho(1+\log\Delta)+\ell)$\\
\bottomrule
\end{tabularx}
\caption{General-graph consequences in the identifier-based \LOCAL\ model.
The additive dependence on $n$ is included.}
\label{tab:general}
\end{table}

The constants in the fixed-$t$ and fixed-$c$ families may depend on those
parameters. Consequently these families have no single smallest member
if only an unspecified constant round budget is given. The asymptotic
expressions have the usual bounded-degree interpretation. Every sufficient
nonintegral list size is rounded up.

\FloatBarrier
\subsection{Prior ingredients and the main questions}
The elementary proof in \cref{sec:hall} uses Hall's theorem~\cite{Hall} and
the probabilistic construction of expanding neighborhoods. This is the
standard lossless-expander counting method; a general existence theorem
appears in Berinde et al.~\cite[Proposition 6]{BGIKS}. We give the full
specialized proof with explicit constants. The sharper ordinary reductions
use Haxell's hypergraph matching condition~\cite{Haxell}; its exact
bipartite formulation also appears in Annamalai~\cite[Theorem 1.1]{Annamalai}.
We include a proof of the needed form.

The proposal-and-acceptance structure is related to the two-round
color-reduction and list-coloring framework of Maus and
Tonoyan~\cite{MT}. The list slack-reduction strategy comes from
Balliu, Kuhn, and Olivetti~\cite[Lemma 4.2]{BKO}, also used in
Balliu, Brandt, Kuhn, and Olivetti~\cite[Appendix D]{BBKO}. We prove the
version needed here and use virtual vertex splitting to obtain only
$O(\rho)$ sequential classes. General-graph reductions use standard
vertex-color reduction~\cite{Linial}, defective coloring~\cite{Kuhn,BEK},
and, for the second row of \cref{tab:general}, the conflict-coloring
theorem of Fraigniaud, Heinrich, and Kosowski~\cite{FHK}.

For comparison, \cite[Lemma 6.1]{BBKO} gives $(2+\eps)\Delta$ colors in
$O(\eps^{-6}\log^{11}\Delta)$ deterministic \CONGEST\ rounds when the
bipartition is supplied. At most $2\Delta-2$ colors have a quite different
behavior: Chang, He, Li, Pettie, and Uitto~\cite[Theorem 2]{CHLPU} prove
$n$-dependent lower bounds even with a supplied bipartition. That theorem
does not give a lower bound for a larger constant times $\Delta$.
We do not establish whether ordinary $O(\Delta)$ coloring is possible in
constant time, or whether arbitrary $100\Delta$-element lists can be
colored in $O(h)$ rounds. The restrictions proved in \cref{sec:barriers}
concern specified constructions and do not settle these questions.

\section{The elementary proof: expansion and Hall twice}\label{sec:hall}

\begin{lemma}[Hall's theorem~\cite{Hall}]\label{lem:hall}
Finite sets $A_i$, indexed by $i\in I$, admit distinct representatives
$a_i\in A_i$ if and only if
$\abs{\bigcup_{i\in J}A_i}\ge\abs J$ for every $J\subseteq I$.
\end{lemma}
\begin{proof}
Necessity is immediate. For sufficiency, induct on $\abs I$.
If a nonempty proper $J\subset I$ has equality, match its sets by
induction and delete their union from the other sets. For any collection
$K\subseteq I\setminus J$, Hall's inequality for $K\cup J$ implies that
at least $\abs K$ colors remain, so induction matches those sets too.
If there is no such tight $J$, choose an element of one set, delete that
set and element, and use induction: every nonempty collection of the
remaining sets originally had at least one unit of slack. The cases
$\abs I\le1$ are immediate.
\end{proof}

\begin{lemma}[Expanding neighborhoods]\label{lem:expander}
For integers $\Delta\ge1$ and $r\ge2$, there are sets
$N_i\subseteq[48r\Delta]$, $i\in[2^r\Delta]$, with
\[
 \abs{N_i}=2r+1,\qquad
 \abs{\bigcup_{i\in I}N_i}\ge(r+1)\abs I
 \quad\text{whenever }\abs I\le\Delta.
\]
Here $[a]=\{1,\ldots,a\}$.
\end{lemma}
\begin{proof}
Choose the $N_i$ independently and uniformly among $(2r+1)$-subsets
of $[Q]$, where $Q=48r\Delta$. A fixed $D$-subset chosen uniformly from
$[Q]$ lies in a fixed $t$-set with probability
\[
 \frac{\binom tD}{\binom QD}
 =\prod_{j=0}^{D-1}\frac{t-j}{Q-j}\le(t/Q)^D
\]
when $t\ge D$; it has probability zero otherwise.
We use $\binom ab\le(ea/b)^b<(3a/b)^b$ for $1\le b\le a$.
For completeness, the binomial theorem gives
$\binom ab(b/a)^b\le(1+b/a)^a\le e^b$.

For a fixed $1\le k\le\Delta$, a union of at most $(r+1)k$ colors is
contained in some set of exactly that size. A union bound gives probability
at most
\begin{align*}
 &\binom{2^r\Delta}{k}\binom{48r\Delta}{(r+1)k}
 \left(\frac{(r+1)k}{48r\Delta}\right)^{(2r+1)k}\\
 &\qquad\le
 \left[9\left(\frac{r+1}{8r}\right)^r
              \left(\frac{k}{\Delta}\right)^{r-1}\right]^k
 \le(81/256)^k<(1/3)^k.
\end{align*}
The second inequality uses $(r+1)/r\le3/2$ and $r\ge2$.
Summing over $k$ gives failure probability less than $1/2$.
The resulting neighborhoods in fact have strictly more than
$(r+1)\abs I$ neighbors for every nonempty eligible $I$.
\end{proof}

\begin{lemma}[One palette reduction]\label{lem:expanderreduction}
An existing proper edge coloring with at most $2^r\Delta$ colors can be
replaced by one with $48r\Delta$ colors in two rounds, for $r\ge2$.
\end{lemma}
\begin{proof}
Associate the fixed neighborhood $N_i$ of \cref{lem:expander} with old
color $i$. At a black vertex, choose an $(r+1)$-subset $S_e\subseteq
N_{c(e)}$ for every incident edge, with the $S_e$ pairwise disjoint.
Such a choice exists: replace each incident edge by $r+1$ clones whose
available colors are $N_{c(e)}$. A set of clones supported on $k$ edges
has size at most $(r+1)k$ and that many available colors by expansion.
Hall matches all clones, giving the required sets. Distinct old colors
on a star are essential here.

In the first round, black vertices send their proposals $S_e$ to white
vertices. Consider any $k$ edges at a white vertex. Their old colors
are distinct and
\begin{equation}\label{eq:deletedneighbors}
 \abs{\bigcup_e S_e}
 \ge\abs{\bigcup_e N_{c(e)}}-\sum_e\abs{N_{c(e)}\setminus S_e}
 \ge(r+1)k-rk=k.
\end{equation}
Hall therefore supplies distinct representatives for all received
proposals. The white vertex assigns these representatives to its edges
and returns them in the second round. They are distinct at white by
construction and at black because the proposals were disjoint.
\end{proof}

This proof explains the choice $2r+1$: expansion by $r+1$ allows black
vertices to allocate $r+1$ colors per edge, and deleting the other $r$
colors per neighborhood still leaves the one unit of expansion needed
at white vertices. The algorithm only needs the neighborhoods and
ordinary bipartite matchings.

\begin{theorem}[Linear palettes in iterated-logarithmic time]\label{thm:linearordinary}
A supplied-bipartite graph of maximum degree at most $\Delta$ has a
uniform deterministic $512\Delta$-edge-coloring algorithm taking $O(h)$
rounds. Identifiers are unnecessary.
\end{theorem}
\begin{proof}
Exchange ports once and color an edge by the ordered pair of its black
and white ports. This is a proper $\Delta^2$-coloring known at both ends.
While the current palette has size $L\Delta$ with integer $L>512$, put
$r=\ceil{\log_2 L}$ and apply \cref{lem:expanderreduction}. The new
overhead is $f(L)=48\ceil{\log_2 L}$. For $r\ge10$,
$48r<2^{r-1}<L$: the first inequality holds at $r=10$ and persists since
$(r+1)/r<2$. Thus the integer overhead strictly decreases to at most $512$.

For $x\ge2$, $f(x)\le96\log_2 x$. For all $x$ above an absolute
threshold, $f(f(x))\le\log_2 x$, since
$96\log_2(96\log_2 x)\le\log_2 x$ eventually. Each pair of reductions
therefore removes at least one logarithmic level until that threshold;
below it there are only an absolute number of decreasing integer steps.
This is $O(1+\logstar\Delta)$ reductions.

The existence argument is uniform: enumerate all possible arrays
$(N_i)$ lexicographically and stop at the first satisfying the finite
expansion tests. Every vertex uses the same parameters and constructs
the same array. Exhaustive search also finds the local matchings. These
computations terminate and consume no communication rounds.
\end{proof}

\subsection{Alternating sides and message sizes}\label{sec:implementation}
After the proposal round of \cref{lem:expanderreduction}, all current
colors are known at the white vertices, and the conceptual coloring is
proper everywhere. One may immediately reverse the roles: white
vertices choose disjoint proposals for the next reduction, and black
vertices accept representatives. A passive vertex need not know old
color indices to compute a matching among the received proposals.
It only needs the guarantee, supplied by properness of the conceptual
old coloring, that those indices are distinct.

Consequently $t\ge1$ consecutive reductions take $t$ proposal rounds
and one final notification. Including the initial port exchange, the
total is $t+2$. A reduction viewed in isolation with agreed input and
output colors still takes two rounds. This also explains the one-round
description in an active/passive formulation of a bipartite problem.

To implement a proposal in \CONGEST, index the possible proposals inside
each old-color neighborhood. Send the old color and this index; the
recipient reconstructs the proposal using the common table. There are
$\binom{2r+1}{r+1}\le2^{2r+1}$ proposals, so the index needs $O(r)$ bits.
Along the reduction sequence $r=O(\log\Delta)$ and every color domain is
polynomial in $\Delta$. All messages use $O(\log\Delta)$ bits, hence
$O(\log n)$ for simple graphs with a known bound $\Delta\le n-1$.
The ordinary constructions also work for loopless bipartite multigraphs with distinct local ports. The relation $\Delta\le n-1$ used for the bandwidth statement is specific to simple graphs. For the sharper reduction of \cref{thm:sharp}, the candidate count satisfies $\log m=O(\log L+r)=O(\log\Delta)$ along the relevant sequences; in \cref{app:constant} it is an absolute constant. The same index encoding therefore applies there as well. This observation addresses bandwidth only; finite-table construction
and storage may be extremely expensive.

\section{Sharper ordinary palette reductions}\label{sec:sharp}

\begin{definition}\label{def:families}
For an input domain $[M]$ and output domain $[Q]$, a palette system is
a nonempty finite family $\F_i$ of nonempty subsets of $[Q]$ for each
$i\in[M]$. It is \emph{good up to degree $\Delta$} if, for every
$I\subseteq[M]$ of size at most $\Delta$:
\begin{enumerate}
\item there are pairwise disjoint choices $S_i\in\F_i$ for $i\in I$;
\item every choice $S_i\in\F_i$ for $i\in I$ admits distinct representatives.
\end{enumerate}
\end{definition}

\begin{lemma}[Palette protocol]\label{lem:protocol}
A good palette system converts an $M$-edge coloring to a $Q$-edge
coloring in two rounds. Its reductions can be pipelined as in
\cref{sec:implementation}.
\end{lemma}
\begin{proof}
Black vertices choose disjoint palettes using the first property and
send them. White vertices use the second property to choose distinct
representatives and return them. This is the proof of
\cref{lem:expanderreduction}, with its two Hall arguments replaced by
the two stipulated properties.
\end{proof}

The following classical matching condition will be used to obtain the
first property. A full proof is given in \cref{app:haxell}.
\begin{theorem}[Haxell's condition~\cite{Haxell,Annamalai}]\label{lem:haxell}
Let $\F_i$, $i\in I$, be nonempty finite families of nonempty sets of
size at most $r$. If, for every nonempty $J\subseteq I$, every set
meeting every palette in all families indexed by $J$ has size greater
than $(2r-1)(\abs J-1)$, then pairwise disjoint choices
$S_i\in\F_i$ exist.
\end{theorem}

\begin{lemma}[Robust Hall estimate]\label{lem:hallbound}
Independently sample $m$ uniform $r$-subsets of $[R\Delta]$ for each of
$L\Delta$ families, where $r\ge2$, $R\Delta\ge r$, and the domain sizes are integers.
If $a=9Lm/R^{r-1}<1$, the probability that the second property of
\cref{def:families} fails is at most $a^2/(1-a)$.
The same bound holds for ranges of independent uniform $r$-tuples.
\end{lemma}
\begin{proof}
A Hall violation gives $s+1$ different families and one palette from
each, all contained in an $s$-set, for $1\le s<\Delta$. To see this,
take a violating collection, let $s$ be its union size, and retain
$s+1$ of the indices. Nonempty palettes exclude $s=0$.
For a fixed $s$ the union bound is
\begin{align*}
 A_s&\le\binom{R\Delta}s\binom{L\Delta}{s+1}
          [m(s/(R\Delta))^r]^{s+1}\\
 &\le\left[9(L/R)m\left(s/(R\Delta)\right)^{r-2}\right]^{s+1}
 \le a^{s+1}.
\end{align*}
In the second line, bound both binomial coefficients by $(3N/k)^k$,
increase the exponent $s$ of $3R\Delta/s$ to $s+1$, and use
$s+1\ge s$. The containment bound applies both to exact subsets and to
tuple ranges. Summing the geometric series proves the claim.
\end{proof}

\begin{theorem}[Sharper reduction]\label{thm:sharp}
For integers $L\ge6$ and $\Delta\ge1$, let $r\ge3$ be the least
integer with $r^r\ge128L^2$. A proper $L\Delta$-edge coloring can be
reduced in two rounds to $8r\Delta$ colors. In particular, the overhead
becomes $O(\log L/\log\log L)$.
\end{theorem}
\begin{proof}
Put
\[
 R=8r,\qquad q=(3/4)^r,\qquad
 m=\ceil{12(L+r)/q}.
\]
For each old color sample $m$ independent uniform exact $r$-subsets of
$[R\Delta]$. First exclude covers of all palettes in any $k$ families
by a set $T$ of size $2r(k-1)$, for $2\le k\le\Delta$.
Writing $x=k/\Delta$, the probability that one sample meets this fixed
$T$ is at most
\[
 g(x)=1-(1-x/4)^r.
\]
Indeed, the exact avoidance probability is
$\prod_{j=0}^{r-1}(1-2r(k-1)/(R\Delta-j))$, and each fraction is at
most $2rk/(R\Delta)=x/4$, since $R\Delta\ge kj$.

The logarithmic elasticity $xg'(x)/g(x)$ decreases on $(0,1]$:
with $y=1-x/4$ it equals
$r/(1+y^{-1}+\cdots+y^{-(r-1)})$. Therefore, with
\[
 \alpha=g'(1)/g(1)=\frac{rq}{3(1-q)},
 \qquad g(x)\le(1-q)x^\alpha.
\]
The last inequality follows by integrating
$d\ln g(x)/d\ln x\ge\alpha$ from $x$ to $1$.
Monotonicity of $z\ln(3R\Delta/z)$ on $0<z\le R\Delta$ gives
\[
 \binom{R\Delta}{2r(k-1)}
 \le\left(\frac{3R\Delta}{2rk}\right)^{2rk}.
\]
Thus the cover-witness probability for size $k$ is at most
\[
 B_k\le[b\,x^{m\alpha-2r-1}]^k\le b^k,
 \quad b=3L\,12^{2r}(1-q)^m,
\]
because $m\alpha\ge4r(L+r)/(1-q)>2r+1$.
Using $mq\ge12(L+r)$, $e>2$, and $L\le2^L$,
\[
 b\le3L\,12^{2r}e^{-12(L+r)}
 \le3L\,2^{-12L}(9/256)^r<1/4.
\]
The sum of the cover probabilities is at most $b^2/(1-b)<1/12$.
Any cover of size at most $(2r-1)(k-1)$ can be padded to size
$2r(k-1)$; all these sizes are below $R\Delta$. Singleton families
have positive hitting number. Hence absence of the witnesses implies
Haxell's condition for every star.

For robust Hall, \cref{lem:hallbound} applies with
$a=9Lm/(8r)^{r-1}$. The ceiling gives $m\le13(L+r)/q$, so
\[
 a\le\frac{156L(L+r)}{(6r)^{r-1}}.
\]
The integer $L$ itself qualifies in the definition of $r$ because
$L^{L-2}\ge6^4\ge128$; thus $r\le L$.
Also $6^{r-1}/r\ge12$ for $r\ge3$, by induction. Consequently
\[
 a\le\frac{312L^2}{12r^r}\le\frac{13}{64}<\frac14.
\]
Both failure sums together are less than $1/6$, proving existence of
a good system. Canonical finite search and \cref{lem:protocol} give
the algorithm. Taking logarithms of $r^r\ge128L^2$ proves the claimed
asymptotic size of the least $r$.
\end{proof}

\begin{corollary}[Fixed-round hierarchy]\label{cor:hierarchy}
For every fixed integer $t\ge1$, a supplied-bipartite graph can be edge
colored using
\[
 O\!\left(\Delta\frac{\log^{(t)}\Delta}{\log^{(t+1)}\Delta}\right)
 \quad\text{colors in }t+2\text{ rounds}.
\]
\end{corollary}
\begin{proof}
Start with the one-round port-pair coloring and apply \cref{thm:sharp}
$t$ times, using \cref{sec:implementation}. For fixed $t$ and sufficiently
large $\Delta$, composition of $O(\log x/\log\log x)$ has the displayed
bound; this follows inductively from
$\log(\log x/\log\log x)=\Theta(\log\log x)$ and the corresponding
next logarithm. Bounded exceptional degrees can retain the original
port-pair coloring, with a larger constant depending on $t$.
\end{proof}

\begin{theorem}[An explicit linear coefficient]\label{thm:smallconstant}
There is a deterministic supplied-bipartite algorithm using
$\lfloor33\Delta/2\rfloor$ colors in $O(h)$ rounds.
\end{theorem}
\begin{proof}
The complete numerical construction in \cref{app:constant} reduces
$512\Delta$ colors through the coefficients
\[
 512\longrightarrow100\longrightarrow20\longrightarrow18
 \longrightarrow17\longrightarrow16.6\longrightarrow16.5
\]
in a constant number of reductions for every $\Delta\ge6000$.
Combine it with \cref{thm:linearordinary}. For $\Delta<6000$, process
the $512\Delta$ old color classes sequentially. Each is a matching;
give each edge a color from $[2\Delta-1]$ avoiding previously assigned
colors. At most $2\Delta-2$ colors are forbidden. A constant number of
rounds per class suffices, and the number of classes is bounded by an
absolute constant in this branch. This gives the assertion for all degrees.
\end{proof}

\section{A direct algorithm for exactly two rounds}\label{sec:tworounds}

The preceding reductions need an initial proper coloring. The following
algorithm works directly from local ports and uses only grouping and ranks.

\begin{theorem}\label{thm:tworounds}
For a supplied-bipartite graph of maximum degree at most $\Delta\ge1$,
there is a deterministic two-round edge-coloring algorithm using
\[
 \Delta(2s+b^2)<3\Delta^{5/3}\quad\text{colors},\qquad
 s=\floor{\Delta^{2/3}},\quad b=\floor{\Delta/(s+1)}.
\]
Both endpoints know the output after the second round, and messages have
$O(\log\Delta)$ bits. The construction also works with parallel edges.
\end{theorem}
\begin{proof}
Exchange ports in round one. For edge $e$, let $p(e)$ be its black port
and $q(e)$ its white port, and put $z(e)=(p(e)+q(e))\bmod\Delta$.
At a black vertex, group edges by their opposite port $q$. Call a group
light if it has at most $s$ edges and heavy otherwise. There are at most
$b$ heavy groups. Give every light edge a zero-based rank $a_B(e)<s$
within its group, in increasing order of $p$. Give a heavy edge a
zero-based rank $b_B(e)$ among \emph{all black-heavy edges with the same
$z$}, again in port order. This rank is less than $b$: for a fixed $z$
and opposite port $q$, there is at most one local port $p\in[0,\Delta-1]$.
Thus each heavy group contributes at most one edge at that $z$.
White vertices compute the symmetric flags and ranks $a_W,b_W$.
These computations use only information received in round one; heavy
ranks do not depend on the other endpoint's classification.

In round two each endpoint sends its light/heavy flag and the applicable
rank. Both endpoints then use the same priority rule:
\[
 c(e)=
 \begin{cases}
 (B,q(e),a_B(e)),&\text{if black calls $e$ light},\\
 (W,p(e),a_W(e)),&\text{otherwise, if white calls $e$ light},\\
 (H,z(e),b_B(e),b_W(e)),&\text{otherwise}.
 \end{cases}
\]
The tags specify disjoint output palettes. Within tag $B$, the pair
$(q,a_B)$ distinguishes edges at black, while the local port $q$
distinguishes them at white. Tag $W$ is symmetric. Within tag $H$,
different $z$ values differ, and at fixed $z$ the appropriate local
heavy rank distinguishes incident edges at each endpoint. Restricting
those ranks to edges heavy at both endpoints preserves distinctness.
Hence the coloring is proper and both endpoints compute the same result.

The tags use at most $\Delta s$, $\Delta s$, and $\Delta b^2$ colors.
Since $s\le\Delta^{2/3}$ and $s+1>\Delta^{2/3}$, we have
$b<\Delta^{1/3}$, proving the strict displayed bound. If $b=0$, the
heavy tag is unused. Each message contains a port, or a flag and a rank,
and therefore has $O(\log\Delta)$ bits.
\end{proof}

\section{Three-round coloring from arbitrary lists}\label{sec:lists}

Ordinary palette reduction does not directly preserve edge lists:
an allowed list may avoid every predetermined candidate palette.
We replace disjoint-choice families by families that are rich inside
every sufficiently large allowed set. We first work in a finite common
universe, and then remove that assumption by a local encoding.

\begin{lemma}[Rich families with robust representatives]\label{lem:rich}
Let $r\ge3$, $C\ge1$, and $r^r\ge180C^2$. For every integer $j\ge1$
with $rj\le C$ there are $j$ families
$\F_{j,1},\ldots,\F_{j,j}$ of nonempty subsets of $[C]$, each of
size at most $r$, such that:
\begin{enumerate}
\item every $rj$-subset of $[C]$ contains a palette from each family;
\item every selection of palettes from distinct families admits distinct
representatives.
\end{enumerate}
\end{lemma}
\begin{proof}
Put $t=rj$ and $m=\ceil{4C(C/t)^r}$. In each family, independently
sample $m$ independent uniform $r$-tuples in $[C]$ and use their ranges
as palettes. A fixed $t$-set contains the range of a fixed tuple with
probability $(t/C)^r$. Thus a union bound for richness failure gives
\[
 j\binom Ct(1-(t/C)^r)^m
 \le C2^C e^{-4C}<C2^{-3C}\le1/8.
\]
For robust Hall, a violation is witnessed by $s+1$ different families
with a palette each inside an $s$-set, for $1\le s<j$. Its probability
is at most
\begin{align*}
 B_s&\le\binom Cs\binom j{s+1}m^{s+1}(s/C)^{r(s+1)}\\
 &\le\left[\frac{9jm s^{r-2}}{C^{r-1}}\right]^{s+1}
 \le a^{s+1},\qquad a=\frac{9jm j^{r-2}}{C^{r-1}}.
\end{align*}
The binomial calculation is the one in \cref{lem:hallbound}.
Since $m\le5C(C/(rj))^r$,
\[
 a\le\frac{45C^2}{r^rj}\le\frac1{4j}\le\frac14.
\]
Summing gives Hall-failure probability at most $1/12$. There are no
Hall witnesses when $j=1$. The total failure probability is below one,
so each required $j$-table exists. It is constructed separately for each
$j$; no union bound over different values of $j$ is needed. The
properties are finite and decidable, hence canonical exhaustive search
gives a common deterministic table.
\end{proof}

\begin{lemma}[Finite-universe list coloring]\label{lem:finiteuniverse}
Under the assumptions $r\ge3$ and $r^r\ge180C^2$, a supplied-bipartite
graph whose lists satisfy
\[
 L_{uv}\subseteq[C],\qquad
 \abs{L_{uv}}\ge r(\deg(u)+\deg(v)-1)
\]
can be list edge colored in three rounds.
\end{lemma}
\begin{proof}
In round one every white vertex sends its degree $j$ and its one-based port
index $i\in[j]$ along each incident edge. The appropriate $j$-table
exists: on every edge at that white vertex, the list-size promise
implies $rj\le C$.

A black vertex of degree $b$ processes its edges in local port order.
For an edge going to white degree $j$, remove the colors in all its
previously chosen palettes. At most $r(b-1)$ colors have been removed,
leaving at least
$r(b+j-1)-r(b-1)=rj$ colors. Richness supplies a palette from
$\F_{j,i}$ contained in the remainder. In round two send that palette
to white. All palettes chosen at a black vertex are disjoint, even if
its incident family indices repeat.

At a white vertex the family indices are its distinct ports in one
fixed $j$-table. Robust Hall gives distinct representatives. Choose
the lexicographically first such assignment and return the colors in
round three. They are allowed by the lists and distinct at both sides.
\end{proof}

\begin{theorem}[Arbitrary labels, three rounds]\label{thm:listprimitive}
Let $\rho$ be defined by \eqref{eq:rho}. If the lists on a
supplied-bipartite graph of maximum degree at most $\Delta$ satisfy
\[
 \abs{L_{uv}}\ge\rho(\deg(u)+\deg(v)-1),
\]
then they can be colored in three deterministic rounds, with no bound
on the global color universe.
\end{theorem}
\begin{proof}
Set $K=2\rho\Delta$ and $C_0=2\rho\Delta^2$. Truncate each edge list
canonically to its first $\min\{\abs{L_e},K\}$ colors; both endpoints
perform the same truncation. Its pointwise promise survives because
$\rho(\deg(u)+\deg(v)-1)<K$.
At a white vertex $w$, form the union $U_w$ of its truncated incident
lists. This union has size at most $C_0$. Encode it injectively in
$[C_0]$ by its order ranks, writing $\phi_w$ for the encoding.
The definition of $\rho$ gives
\[
 \rho^\rho\ge720\rho^2\Delta^4=180C_0^2.
\]
Thus \cref{lem:rich} supplies all necessary tables on $[C_0]$.

In round one $w$ sends the ordered union $U_w$, its degree, and its
port along each incident edge. A black vertex now knows the encoding
and inverse for each of its edges. It chooses palettes greedily as in
\cref{lem:finiteuniverse}, but removes colors using their \emph{actual
labels}. For the current edge, encode an available $\rho j$-set using
the appropriate $\phi_w$, choose a palette inside that image by
richness, and decode it to actual labels. Send the actual palette in
round two. No unused coordinate is chosen because the palette is
contained in an image of actual allowed colors.

In round three each white vertex encodes its received palettes using
its single map $\phi_w$, chooses distinct representatives using its
table, and sends back their actual labels. Injectivity preserves
distinctness at white. At black, disjointness was enforced on actual
labels across all incident edges, including those using different
encodings. This proves correctness and the exact round count.
\end{proof}

One shared encoding per white star is crucial. Independent rank
encodings of individual edge lists would not preserve whether two
incident actual colors are equal. The size of the messages here is
unrestricted; this is a \LOCAL\ result.

\section{List slack, batching, and partial coloring}\label{sec:slack}

\subsection{Batching with constant additive slack}
\begin{proposition}\label{prop:quadraticlists}
Lists of size at least $\Delta^2$ can be edge colored in two rounds
on a supplied-bipartite graph.
\end{proposition}
\begin{proof}
Each black vertex chooses disjoint $\Delta$-element sublists for its
edges greedily. At most $\Delta(\Delta-1)$ colors are excluded before
any choice, so this is possible. Send them in round one. A white
vertex receives at most $\Delta$ sets, all of size $\Delta$, so every
nonempty subcollection has union size at least its cardinality.
Hall gives distinct representatives; return them in round two.
\end{proof}

\begin{theorem}[Batching tradeoff]\label{thm:batching}
For $1\le k\le\Delta$, lists of size
\begin{equation}\label{eq:batchsize}
 2\Delta-2+\rho(2k-1)
\end{equation}
suffice in $O(h+\Delta/k)$ rounds on supplied-bipartite graphs.
More generally, lists of size $\ceil{(2+\eps)\Delta}$ suffice in
$O(h+1+\rho/\eps)$ rounds for every $\eps>0$.
\end{theorem}
\begin{proof}
Compute a proper $512\Delta$ auxiliary edge coloring using
\cref{thm:linearordinary}. Group its colors into consecutive batches
of at most $k$. A batch induces maximum degree at most $k$.
Process batches sequentially, removing colors used on previously
colored adjacent edges. An edge has at most $2\Delta-2$ adjacent
edges in total, so \eqref{eq:batchsize} leaves at least
$\rho(2k-1)$ colors. Apply \cref{thm:listprimitive} to the batch,
keeping the ambient parameters $\Delta,\rho$ fixed.
Its pointwise list requirement is no larger than that remainder.
Between calls, endpoints exchange the sets of colors already used
at their vertices, so both endpoints agree on the residual list.
This takes a constant number of rounds per batch. Properness within
a batch is supplied by the primitive, and across batches by deleting
previously used colors. There are $O(\Delta/k)$ batches.

For the second assertion put $x=\eps\Delta/(2\rho)$. If $x\ge1$,
take $k=\min\{\Delta,\floor x\}$. At least $\eps\Delta+2$
colors remain on each edge, which is enough since
$\rho(2k-1)\le2\rho k\le\eps\Delta$.
When $1\le x<\Delta$, $\floor x\ge x/2$, hence the number of batches
is $O(1+\rho/\eps)$; the same bound holds when $x\ge\Delta$.
If $x<1$, process matching classes directly, choosing any remaining
color. There are $O(\Delta)=O(\rho/\eps)$ classes.
\end{proof}

The simpler quadratic-list primitive can similarly replace
$\rho(2k-1)$ in \eqref{eq:batchsize} by $k^2$. Thus, if useful for
particular finite parameters, $2\Delta-2+\min\{k^2,\rho(2k-1)\}$
also suffices with the same asymptotic time. Setting $\eps=98$ proves
the $100\Delta$-list bound in \cref{tab:bipartite}. The choices of $k$
described after that table give its other consequences.

\subsection{Degree-plus-one lists by halving residual degrees}
We give a full slack-reduction proof in the style of
\cite[Lemma 4.2]{BKO}. The improvement in the number of sequential
classes comes from the ordinary coloring theorem and virtual vertex
splitting.

\begin{theorem}[Partial coloring and degree-plus-one lists]\label{thm:partial}
On a supplied-bipartite graph, suppose $\abs{L_e}\ge d_G(e)+1$.
There is a phase taking $O(\rho)$ rounds which preserves this invariant
and reduces the maximum residual edge degree by at least a factor of
two, unless all edges are colored. Consequently all edges can be
colored in $O(\rho(1+\log\Delta))$ rounds. More generally $t$ phases
cost $O(\rho t)$ rounds and leave maximum edge degree at most
$(2\Delta-2)/2^t$.
\end{theorem}
\begin{proof}
Let $H$ be the residual graph at the start of a phase and freeze
$D_e=d_H(e)$. Set $p=4\rho$. At every vertex, group its incident
residual edges in local port order into blocks of at most $p$, and
replace each block by a virtual vertex with the same bipartition side.
The resulting virtual graph has maximum degree at most $p$ and the
same edges. One physical vertex simulates its blocks, and each virtual
edge uses its corresponding physical edge. Thus one virtual round
costs one physical round.

Apply \cref{thm:linearordinary} with degree bound $p$, obtaining at
most $512p$ auxiliary edge colors in $O(1+\logstar p)$ rounds. Use
these colors as classes on the original graph. A class contains at
most one edge from each block. If $e=uv$ belongs to that class, its
edge degree within the class is at most
\begin{align*}
 \ceil{\deg_H(u)/p}-1+\ceil{\deg_H(v)/p}-1
 &=\floor{(\deg_H(u)-1)/p}+\floor{(\deg_H(v)-1)/p}\\
 &\le D_e/p.
\end{align*}

Process the classes in order. Delete colors already used by colored
adjacent edges, and declare an edge active if its remaining list has
more than $D_e/2$ colors. Endpoints synchronize these lists and the
active status in a constant number of rounds. In the active subgraph,
the current edge degree $D'_e$ is at most $D_e/(4\rho)$, hence
$\abs{L_e}>2\rho D'_e$. If $D'_e\ge1$, this implies
$\abs{L_e}\ge\rho(D'_e+1)$, so \cref{thm:listprimitive} colors the
active nonisolated edges in three rounds. Isolated active edges may
have shorter lists; color them directly from any available color.
They are separate components and can be handled concurrently.

An edge left uncolored was inactive when its class was processed.
Its list started with at least $D_e+1$ colors and then had at most
$D_e/2$. Therefore at least $D_e/2+1$ colors had been deleted, which
requires at least that many adjacent edges to have been colored.
Its final residual degree is at most $D_e/2-1$.
At every stage, coloring an adjacent edge removes one adjacency and
deletes at most one list color, preserving the degree-plus-one
invariant. In particular, an isolated residual edge is active and
will be colored. The maximum residual edge degree halves, and after
$O(1+\log\Delta)$ phases no edge remains.

Each of $512p=O(\rho)$ classes costs constant time, and
$1+\logstar p=O(\rho)$. The phase time and the partial-coloring claim
follow. All degree and list information used in a class concerns the
current subgraph; the ambient parameters $\Delta,\rho$ remain valid.
\end{proof}


\section{General graphs}
\label{sec:general}

Whenever an algorithm is applied to a subgraph, its incident edges are renumbered locally from zero in their inherited port order; the algorithms use these subgraph ports.

We now consider simple graphs with identifiers, a known maximum-degree
bound $\Delta$, and arbitrary finite edge lists known at both endpoints.
Write $h=1+\log^*\Delta$ and $\ell=1+\log^*n$.  The parameter
$\rho=\rho(\Delta)$ of Theorem~\ref{thm:listprimitive} is kept fixed when
we pass to subgraphs.  We first explain the ordinary-coloring reduction,
then give list-preserving reductions.  In particular, multiplying an
ordinary color by a refinement color does not preserve arbitrary lists.

\subsection{Ordinary coloring by splitting vertices}

\begin{theorem}\label{thm:generalordinary}
A deterministic LOCAL algorithm properly edge-colors a graph of maximum
degree at most $\Delta$ with $O(\Delta)$ colors in $O(\ell)$ rounds.
More precisely, a supplied-bipartite $q(\Delta)$-edge-coloring algorithm
taking $T(\Delta)$ rounds yields a general-graph algorithm using
$3q(\Delta)$ colors in $O(T(\Delta)+\ell)$ rounds.
\end{theorem}

\begin{proof}
Orient every edge from its smaller-identifier endpoint to its
larger-identifier endpoint.  Replace each vertex $v$ by a black copy
$v^+$ and a white copy $v^-$, and replace $u\to v$ by $u^+v^-$.  The
resulting virtual graph has maximum degree at most $\Delta$.  A physical
vertex simulates its two copies, so one virtual communication round costs
one physical round.  Apply the supplied-bipartite algorithm and project
the resulting colors onto the original edges.

For every projected color, each vertex has at most one incoming and at
most one outgoing edge of that color.  Its monochromatic components are
directed paths: a cycle would be directed, contradicting the strictly
increasing identifiers along directed edges.  A standard deterministic
constant-degree coloring algorithm properly $3$-edge-colors all these
paths in parallel in $O(\ell)$ rounds. Explicitly, host the process for
each path edge at its tail, using its tail identifier as its name. These
names are distinct within each path. Apply the proper vertex-coloring
algorithm proved in Appendix~\ref{app:defective} to the line graph of
each path, whose maximum degree is two, to obtain a constant-size
proper palette. Process those color classes sequentially and recolor
into $[3]$, avoiding the at most two previously assigned neighbor
colors. The simulation has constant cost per line-graph round.  Each edge belongs to one path instance,
and all instances can be simulated concurrently.  Output the pair
consisting of the projected color and the path color.  Edges with
different first coordinates differ, and adjacent edges with the same
first coordinate differ in the second coordinate.  The first claim now
follows from Theorem~\ref{thm:linearordinary}, since $h=O(\ell)$.
\end{proof}

\subsection{Two standard vertex-coloring tools}

The list reductions use the following published distributed
vertex-coloring tools. We state their needed quantitative forms here and give complete proofs in Appendix~\ref{app:defective}.

\begin{lemma}[Standard color reduction and defective coloring]
\label{lem:standarddefect}
There is an absolute constant $A$ with the following properties.
\begin{enumerate}
\item A graph of maximum degree $D\ge1$ with a proper $M$-vertex-coloring
can obtain a proper $A D^2$-vertex-coloring in $O(1+\log^*M)$ rounds.
In particular, identifiers of polynomial range in $n$ give such a
coloring in $O(\ell)$ rounds.
\item For every integer $b$ with $1\le b\le D$, the same input admits
a vertex-coloring with $O((D/b)^2\log(2D/b))$ colors for which every
monochromatic subgraph has maximum degree at most $b$, in
$O(1+\log^*M)$ rounds.  The implied constant in the running time is
independent of $b$.
\end{enumerate}
\end{lemma}

The first part is Linial's color reduction~\cite{Linial}. For the
second, our proof follows the defective-color-reduction method of
Kuhn~\cite{Kuhn}; the stronger published palette bound omits the
logarithmic factor. See also~\cite{BEK}. The slightly weaker palette
here has no effect on any of our edge-coloring bounds. Starting from a
proper $O(\Delta^2)$-coloring, both tools take $O(h)$ rounds on every
subgraph. We retain that original proper coloring throughout.

\begin{lemma}[Binary decomposition]\label{lem:binarydecomposition}
Suppose vertices have labels in $[m]$, not necessarily forming a proper
coloring.  The edges with distinct endpoint labels can be partitioned
into at most $\lceil\log_2m\rceil$ supplied-bipartite subgraphs in one
round.  Edges whose endpoint labels agree are omitted.
\end{lemma}

\begin{proof}
Encode the labels by binary strings of a common length.  Assign an edge
with different endpoint labels to the first bit position in which its
labels differ.  At that position the bit supplies the black/white
bipartition.  Every eligible edge has exactly one such first position.
Exchanging labels lets both endpoints determine the edge's part.
\end{proof}

\subsection{Degree-plus-one lists}

\begin{theorem}\label{thm:generaldegreeone}
If every edge $uv$ has a list of size at least
$\deg(u)+\deg(v)-1$, then the graph can be list edge-colored in
\[
 O\bigl(\rho(1+\log\Delta)+\ell\bigr)
 =O\left(\frac{\log^2\Delta}{\log\log\Delta}+\ell\right)
\]
rounds, with bounded-logarithm conventions for small $\Delta$.
\end{theorem}

\begin{proof}
Initially obtain a proper $O(\Delta^2)$-vertex-coloring in $O(\ell)$
rounds.  At all times delete the colors of already colored adjacent
edges from every uncolored edge's list.  If $H$ is the current residual
graph, the invariant
\[
 |L_e|\ge \deg_H(u)+\deg_H(v)-1
\]
holds: coloring one adjacent edge removes one adjacency and deletes at
most one available color.

Consider a stage whose residual maximum-degree bound is $D$.  For
$D\ge2$, Lemma~\ref{lem:standarddefect} gives a vertex-coloring with
at most $c$ colors and monochromatic degree at most $\lfloor D/2\rfloor$
in $O(h)$ rounds, where $c$ is an absolute constant.  Process each pair
of distinct vertex colors in turn.  Its uncolored edges form a graph
with a supplied bipartition and inherit the degree-plus-one list
invariant.  Apply $t=\lceil\log_2(16c)\rceil$ phases of the partial
coloring algorithm from Theorem~\ref{thm:partial}.  This costs $O(\rho)$
rounds and leaves maximum edge degree at most
\[
 \frac{2D-2}{2^t}\le\frac{D}{8c}.
\]
In any graph with maximum edge degree $B$, every nonisolated vertex has
degree at most $B+1$, because any edge incident to it has at least
$\deg(v)-1$ adjacent edges.  Thus, when $D\ge8c$, each processed
pair leaves maximum vertex degree at most $D/(4c)$.  Later pair calls
can only remove further edges.  Every vertex belongs to at most $c-1$
pairs, so all the remaining bichromatic edges together have maximum
degree at most $D/4$.  The unprocessed monochromatic edges have maximum
degree at most $D/2$.  The residual maximum degree is therefore at most
$3D/4$.

Update the common bound to $D\leftarrow\lfloor3D/4\rfloor$, which is valid since actual degrees are integers, and repeat until this bound is less than $8c$.  There are
$O(1+\log\Delta)$ stages, each costing $O(h+\rho)=O(\rho)$ rounds.
For the final graph, reduce the retained proper vertex-coloring to a
constant-size proper vertex-coloring in $O(h)$ rounds.  This gives a
constant-size proper auxiliary edge-coloring: orient edges by vertex
color, and label an edge by its two endpoint colors and its port
numbers at both endpoints.  Edges sharing an endpoint have
different labels.  Greedily process its constantly many matching
classes.  The invariant supplies an available color for every edge.
All chosen colors belong to the original lists.

This reduction is the list-preserving defective-coloring decomposition
used by Balliu, Brandt, Kuhn, and Olivetti~\cite[Theorem~D.4]{BBKO},
with the stronger partial-coloring routine proved here.  The explicit
$O(h)$ decomposition cost is absorbed by $\rho$.
\end{proof}

\subsection{Nearly \texorpdfstring{$2\Delta$}{2 Delta} colors in logarithmic time}

\begin{theorem}\label{thm:generalbatching}
For $1\le k\le\Delta$, lists of size
\[
 L(k)=2\Delta-2+\rho(2k-1)
\]
suffice in
\[
 O\left(\ell+\frac{\Delta}{k}
       +h\log\frac{2\Delta}{k}+\log(2\Delta)\right)
\]
rounds on general graphs.
Consequently, lists of size
\[
 2\Delta+O\!\left(\frac{\Delta}{\log\log\Delta}\right)
\]
suffice in $O(\log(2\Delta)+\ell)$ rounds.  In particular, for every
fixed $\varepsilon>0$, lists of size $\lceil(2+\varepsilon)\Delta\rceil$
suffice within this time bound.  The same assertions hold for ordinary
edge-coloring by using a common list on all edges.
\end{theorem}

\begin{proof}
Obtain and retain a proper $M=O(\Delta^2)$-vertex-coloring in
$O(\ell)$ rounds.  At all times an uncolored edge has at least
$\rho(2k-1)$ available colors, since at most $2\Delta-2$ adjacent
edges can ever be colored.

As long as the current maximum-degree bound $D$ exceeds $k$, compute
the constant-color defective vertex-coloring with monochromatic
degree at most $\lfloor D/2\rfloor$ given by
Lemma~\ref{lem:standarddefect}.  Its cost is $O(h)$.
Process the constantly many pairs of distinct vertex colors
sequentially.  For a pair, compute a proper $O(D)$ auxiliary
edge-coloring by Theorem~\ref{thm:linearordinary}, and group these
auxiliary colors into batches of at most $k$.  Theorem~\ref{thm:listprimitive}
colors each batch in constant time, since its maximum degree is at
most $k$.  Endpoints exchange colors used at their vertices between
successive calls to keep their residual lists identical.  A whole
stage costs $O(h+D/k)$ rounds and leaves only monochromatic edges,
hence maximum degree at most $\lfloor D/2\rfloor$.

There are $O(1+\log(\Delta/k))$ such stages, and the sum of their
degree bounds is $O(\Delta)$.  Their total cost is
\[
 O\left(h\log\frac{2\Delta}{k}+\frac{\Delta}{k}\right).
\]
Once the residual degree is at most $k$, apply
Lemma~\ref{lem:binarydecomposition} to the retained proper
$M$-vertex-coloring.  It partitions \emph{all} residual edges into
$O(\log(2\Delta))$ supplied-bipartite subgraphs, each of maximum degree
at most $k$.  Color these sequentially with the three-round primitive.
This proves the stated general bound without any further degree
decomposition.

For $\Delta\ge2$, set
$k=\lceil\Delta/\log_2\Delta\rceil$.  Then $1\le k\le\Delta$,
$\Delta/k=O(\log\Delta)$, and
$h\log(2\Delta/k)=O(h\log\log\Delta)=O(\log\Delta)$.
Moreover,
\[
 L(k)=2\Delta-2+
 \rho\left(2\left\lceil\frac{\Delta}{\log_2\Delta}\right\rceil-1\right)
 =2\Delta+O\!\left(\frac{\Delta}{\log\log\Delta}\right).
\]
The degree-one case is immediate.  For fixed $\varepsilon$, the
displayed slack is at most $\varepsilon\Delta$ for sufficiently large
$\Delta$.  The finitely many smaller degree bounds are handled by a
constant-degree greedy schedule after color reduction; its constant
may depend on $\varepsilon$.
\end{proof}

\subsection{Short-time general-graph list tradeoffs}

For completeness we give a family of sufficient list sizes for a total
$O(\ell)$-round algorithm.  It uses the following consequence of conflict coloring, whose complete proof is included in Appendix~\ref{app:fhk}.

\begin{lemma}[A consequence of local conflict coloring]
\label{lem:fhkfinish}
There is an absolute constant $B$ such that a graph of maximum degree
$D\ge1$, supplied with a proper $s$-vertex-coloring, can be list-colored
in $O(1)$ rounds if every list has at least
\[
 \max\{s,D,B D^2\log(2D)\}
\]
colors.  The labels of these colors may be arbitrary finite strings.
\end{lemma}

This is the list-coloring specialization of the method of
Fraigniaud, Heinrich, and Kosowski~\cite[Lemma 3.3]{FHK}.
Appendix~\ref{app:fhk} proves it directly, using two transformations of
conflict-coloring instances and a finite lookup table. Its local
relabelings preserve conflicts, so the original labels need no common
bounded universe. Degree zero is immediate. At degree one, components
are isolated vertices or single edges, and one exchange of lists and
proper auxiliary colors permits the smaller-color endpoint to choose
first and the other to avoid its choice. When applied to a line graph,
each round has constant simulation cost in the original graph: host
each edge process at a deterministically selected endpoint, communicate
through the shared endpoint of adjacent edges, and finally notify both
endpoints of the output.

\begin{theorem}\label{thm:generalshorttime}
For every fixed real $a>0$, there is a constant $C_a$ such that lists of
size at least
\[
 \left\lceil C_a\,
   \frac{\Delta^2\log(2\Delta)}{2^{a(1+\log^*\Delta)}}\right\rceil
\]
admit a deterministic list edge-coloring in $O_a(\ell)$ rounds on
general graphs.  Thus these are a family of sufficient upper bounds,
with constants depending on $a$.
\end{theorem}

\begin{proof}
Initially obtain a proper $O(\Delta^2)$-vertex-coloring and a proper
$s=O(\Delta)$ auxiliary edge-coloring in $O(\ell)$ rounds, using
Lemma~\ref{lem:standarddefect} and Theorem~\ref{thm:generalordinary}.
Put $p=\lceil2^{ah/2}\rceil$.
For each fixed $a$, all sufficiently large $\Delta$ satisfy
$p\le\Delta/2$.  Put $b=\lfloor\Delta/p\rfloor$ and compute a
$b$-defective vertex-coloring with $O((\Delta/b)^2\log(2\Delta/b))=O(p^2\log(2p))$ colors
in $O(h)$ rounds by Lemma~\ref{lem:standarddefect}.

Apply Lemma~\ref{lem:binarydecomposition} to these defective colors.
It partitions the bichromatic edges into $O(\log p)=O_a(h)$
supplied-bipartite subgraphs, each of maximum degree at most $\Delta$.
Color these sequentially with Theorem~\ref{thm:listprimitive}.
This is feasible for a sufficiently large $C_a$, since
\[
 \frac{\Delta^2\log(2\Delta)}{2^{ah}}
   =\Delta^{2-o(1)}\gg \rho\Delta+\Delta,
\]
and at most $2\Delta-2$ colors can be removed from any list throughout
the entire process.  The residual graph has maximum degree at most
$b$, consisting only of monochromatic edges of the defective
vertex-coloring.

Its line graph has maximum degree at most $2b-2$ and inherits the
proper $s$-coloring of its vertices from the auxiliary edge-coloring.
Furthermore,
\[
 b^2\log(2b)
 \le\frac{\Delta^2\log(2\Delta)}{p^2}
 \le\frac{\Delta^2\log(2\Delta)}{2^{ah}}.
\]
For sufficiently large $C_a$, the remaining lists therefore satisfy
Lemma~\ref{lem:fhkfinish}, including its $s=O(\Delta)$ requirement.
Finish in constant time using that lemma on the line graph.  Every
color used is an actual member of the edge's original list.

The extra time beyond initialization is $O_a(h)$, and $h=O(\ell)$.
All degree bounds excluded by ``sufficiently large'' form a finite set
depending only on $a$.  Increase $C_a$ so that their lists have at
least $2\Delta-1$ elements and use a constant-degree greedy schedule
after the initial symmetry breaking.  This covers every $\Delta\ge1$.
\end{proof}

\section{Lower bounds and limits of the constructions}\label{sec:barriers}

\subsection{The exact one-round threshold}
\begin{theorem}\label{thm:oneround}
Fix $\Delta\ge2$ and a palette size $q$ independent of $n$. Under
our initialization and two-endpoint output conventions, any deterministic
one-round algorithm valid on all sufficiently large $\Delta$-regular
supplied-bipartite graphs requires $q\ge\Delta^2$, even with identifiers
and knowledge of $n$. The bound also holds for randomized algorithms with
any fixed positive global success probability on all such graphs, if
disconnected inputs are allowed. The port-pair algorithm attains $\Delta^2$.
\end{theorem}
\begin{proof}[Deterministic case]
Write $d=\Delta$ and use $H=K_{d,d}$ as a fixed core with a specified
bipartition. Let $P$ be a pool of $K$ identifiers, where $K$ is a multiple
of $2d$, $K\ge3d+1$, and $K$ is large enough for the finite Ramsey
theorem with $q^{d^2}$ edge colors and a homogeneous set of $d+1$
vertices. There are finitely many core tests: all injective assignments
of pool identifiers to $H$ and all port numberings, in total
\[
 N=(K)_{2d}(d!)^{2d}.
\]
Fix the global parameter $n$ before considering these tests. For
sufficiently large $n$ divisible by $2d$, each test can be padded to
that $n$ with disjoint regular components. Padding is invisible in one round.

Suppose a deterministic protocol succeeds on every core test. For an
edge whose black and white endpoint identifiers are $x,y$ and whose
ports are $p,s$, agreement forces the common output to be a function
$f_n(x,y,p,s)$ of just these data. To prove this, hold the white
endpoint's star fixed while varying the identifiers and opposite ports
at the other neighbors of the black endpoint. These changes cannot affect
the white endpoint's first-round view: first-round messages depend only
on the sender's own initial information. For any two choices of the
black star, choose one white star whose private identifiers avoid both
choices. At most $3d-1$ identifiers are needed. These compatible stars
are realizable in $K_{d,d}$; complete unused vertices with distinct
identifiers. The edges between private vertices add no initial
information. Comparing through this common white star removes dependence
on the black private star. The symmetric comparison removes the white
private star, proving the claim without violating identifier uniqueness.

Color every unordered pool pair $x<y$ by the matrix
$(f_n(x,y,p,s))_{p,s\in[d]}$. There are at most $q^{d^2}$ matrices.
Ramsey gives $d+1$ identifiers on whose pairs the matrix is constant;
denote it by $(a_{p,s})$. At a black center with the smallest identifier
in this set, different local ports $p,p'$ may have arbitrary opposite
ports $s,s'$, forcing $a_{p,s}\ne a_{p',s'}$ by properness.
At a white center with the largest identifier, different local ports
$s,s'$ may have the same opposite port $p$, forcing
$a_{p,s}\ne a_{p,s'}$. All these stars occur among the tests.
Therefore all $d^2$ entries differ, and $q\ge d^2$.

Only the elementary finite form of Ramsey's theorem is used. Its
existence follows by successively choosing a vertex and retaining a
largest color class among its edges to the unused vertices. For $c$
colors, a sufficiently large finite starting set permits $c(k-1)+1$
such choices, each with a fixed color toward all later choices. At
least $k$ chosen vertices have the same fixed color and form a
monochromatic clique. Thus $K$ depends only on $d,q$, not on $n$ or
on the algorithm.
\end{proof}

\begin{proof}[Randomized case]
The preceding proof shows that every deterministic protocol with
$q<d^2$ fails on at least one of the $N$ core tests, for each fixed $n$.
Take $n=tK$ and partition its identifiers into $t$ disjoint pools of
size $K$. In each pool independently choose a uniformly random core
test and fill its unused identifiers by a fixed collection of copies
of $K_{d,d}$. This defines a distribution on valid $n$-vertex inputs.

Fix all algorithmic randomness first: one public tape and a private
tape for each identifier. In each pool the resulting deterministic
algorithm fails at least one of its $N$ tests. The preceding argument
applies to any ordered identifier pool. Because input choices in
different pools are independent and outputs in a component depend only
on that component and the fixed global parameters,
\[
 \Pr_{\mathrm{input}}[\text{global success}\mid\text{all tapes}]
 \le(1-1/N)^t.
\]
Average over the tapes. Some fixed input has algorithmic success
probability at most $(1-1/N)^t$, tending to zero. This excludes any
fixed positive global success guarantee. The argument uses independence
of the input choices after conditioning on all tapes, not independence
of component failures under public randomness.
\end{proof}

The quantifiers fix $d,q$ and allow arbitrarily large $n$; the Ramsey
threshold may be enormous. A separate promise $n=\poly(d)$ is not
covered. Nor does the proof apply unchanged if opposite ports or
neighbor identifiers are supplied initially, or if only one endpoint
must output. The randomized amplification uses disjoint unions. The
theorem gives no superconstant round lower bound for a linear palette.

\subsection{The small-palette obstruction}
Chang et al.~\cite[Theorem 2]{CHLPU} prove lower bounds
$\Omega(\log_\Delta n)$ deterministically and
$\Omega(\log_\Delta\log n)$ randomly for $(2\Delta-2)$-edge coloring,
even with a supplied vertex bipartition. Its reduction explains the
threshold: divide the colors into two groups of $\Delta-1$ and orient
their edges in opposite directions across the bipartition. Every
degree-$\Delta$ vertex sees both groups, so the orientation has no sink.
The published sinkless-orientation lower bound therefore applies, on
high-girth regular instances. The relevant neighborhoods are trees;
a finite $\Delta$-regular tree is not needed.

For any fixed $\eps>0$, choose $\Delta$ large enough that
$(2-\eps)\Delta\le2\Delta-2$, and then let $n$ grow. No round bound
depending only on $\Delta$ can provide this palette uniformly over
$n$. This leaves $2\Delta-1$ and $2\Delta$ untouched and gives no
superlinear two-round palette lower bound. Without a supplied
bipartition, the usual $\Omega(\logstar n)$ symmetry-breaking
obstruction already occurs on cycles for a fixed finite
palette~\cite{Linial}; hence the additive $n$-term in
\cref{tab:general}.

\subsection{A limitation of fixed-fraction expansion}
\begin{proposition}\label{prop:expanderbarrier}
Let $M\ge\Delta$ left vertices have neighborhoods of common size
$a\ge1$ in a $Q$-element right universe. Suppose
\[
 \abs{N(I)}\ge\beta a\abs I\quad(\abs I\le\Delta),
 \qquad 0<\beta\le1,\quad\beta\Delta\ge2.
\]
Then $Q>(e\beta/2)\Delta\ln(M/\Delta)$. In particular constant-fraction
expansion requires $Q=\Omega(\Delta\log(M/\Delta))$.
\end{proposition}
\begin{proof}
Put $t=\ceil{\beta a\Delta}-1$. Expansion gives $Q\ge\beta a\Delta$,
so $a\le t<Q$. Every right-side $t$-set contains the full neighborhoods
of at most $\Delta-1$ left vertices. Choose such a $t$-set $T$ uniformly.
For any particular left vertex,
\[
 \Pr[N_i\subseteq T]
 =\frac{\binom ta}{\binom Qa}
 \ge\left(\frac{t-a+1}{Q}\right)^a
 \ge\left(\frac{\beta a\Delta}{2Q}\right)^a,
\]
using $\beta\Delta\ge2$. Linearity of expectation gives
$M(\beta a\Delta/(2Q))^a\le\Delta-1<\Delta$.
Taking logarithms and maximizing $(\ln z)/z$ over positive $z$ gives
\[
 \ln(M/\Delta)<a\ln\frac{2Q}{\beta a\Delta}
 \le\frac{2Q}{e\beta\Delta}.
\]
\end{proof}

This applies to \cref{lem:expander}, whose expansion fraction exceeds
$1/2$, for $\Delta\ge4$. Its universe size is thus optimal in order
under that expansion requirement. The sharper families of
\cref{thm:sharp} need not arise from such neighborhoods.

\subsection{A limitation of universal list richness}
\begin{proposition}\label{prop:richnessbarrier}
Suppose $C\ge2r\Delta$ and there are $\Delta$ families of nonempty
palettes in $[C]$, each palette of size at most $r$. Assume every
$r\Delta$-set contains a palette from each family, and every choice
of one palette from each family has distinct representatives. Then
\[
 (4r)^r\ge\frac{\Delta+1}{\ln(2\Delta)}.
\]
In particular $r=\Omega(\log\Delta/\log\log\Delta)$.
\end{proposition}
\begin{proof}
Fix a set $W$ of $2r\Delta$ colors. An inclusion-maximal disjoint
packing of palettes from one family inside $W$ contains at least
$\Delta+1$ palettes. Otherwise its union has size at most $r\Delta$;
its complement contains an $r\Delta$-set and hence, by richness,
another disjoint palette, a contradiction.

Include each color of $W$ independently in a random set $T$ with
probability $1/(4r)$. Its expected size is $\Delta/2$, so
$\Pr[\abs T\ge\Delta]\le1/2$. Use $\Delta+1$ disjoint packed palettes
from each family. Each lies inside $T$ with probability at least
$(4r)^{-r}$, independently within that packing. The probability that
$T$ contains no palette from a given family is at most
$\exp(- (\Delta+1)/(4r)^r)$.
If the claimed inequality failed, the union bound
\[
 \Pr[\abs T\ge\Delta\text{ or some family has no palette in }T]
 \le\frac12+\Delta\exp\left(-\frac{\Delta+1}{(4r)^r}\right)<1
\]
would give a set of fewer than $\Delta$ colors containing a palette
from every family, contradicting distinct representatives.
Taking logarithms yields
$r\ln(4r)\ge\ln(\Delta+1)-\ln\ln(2\Delta)$ and hence the assertion,
indeed with leading lower bound $(1-o(1))\ln\Delta/\ln\ln\Delta$.
\end{proof}

The table in \cref{thm:listprimitive} meets these hypotheses at white
degree $\Delta$, with $C=2\rho\Delta^2$. Merely improving its probability
estimates cannot make its palette-size parameter constant. This limits
the specified universal richness and robustness requirements. Families
adapted to more communicated list information, weaker compatibility
conditions, or an iterative list-preserving reduction may avoid it.

\section{What the results leave open}
The elementary algorithm isolates a useful asymmetry: one endpoint
coordinates proposals for its incident edges, and the other coordinates
acceptance. Expansion makes the independent local decisions compatible.
For vertex coloring there is no analogous division of constraints into
the two endpoints of a variable; transferring this mechanism needs
further work.

The main quantitative questions are constant-round ordinary $O(\Delta)$
coloring with a supplied bipartition, $O(h)$ rounds for arbitrary linear
lists, and smaller final palette coefficients in $O(h)$ rounds.
The exact two-round gap remains between a linear lower threshold and
the $O(\Delta^{5/3})$ upper bound. The construction-specific barriers
above do not resolve these distributed questions.

There is also a representational question. A palette family can encode
an enormous collection of local proposals while admitting a short
expansion or matching proof. This fits the round-elimination
viewpoint~\cite{Brandt}, where the active/passive proposal step is one
round. A symbolic treatment retaining useful relaxations of such families
could help explain algorithms without enumerating their full finite
descriptions. No additional algorithmic bound is asserted here.

\appendix



\section{A proof of Haxell's matching condition}\label{app:haxell}
This is the bipartite hypergraph matching condition of Haxell
\cite{Haxell,Annamalai}, with a hyperedge consisting of one family index and
at most $r$ universe elements.  We include its alternating-sequence
proof to keep all matching arguments self-contained.

\begin{proof}[Proof of Theorem~\ref{lem:haxell}]
Induct on $n=|I|$, with $n=0$ immediate.  By induction, fix a matching $M$ choosing disjoint
palettes for all families except one root $i_0$.  A matching element
retains its family label.  We show how to augment $M$.

Maintain a sequence of candidate palettes $a_1,\ldots,a_t$ and nonempty
sets $B_1,\ldots,B_t$ of matching elements, with these invariants:
\begin{enumerate}
\item The candidate palettes are pairwise disjoint.
\item $B_j$ consists of all elements of $M$ intersecting $a_j$, and
      the sets $B_j$ are pairwise disjoint as sets of matching elements.
\item The family of $a_j$ is the root or the family of a matching
      element in $B_1\cup\cdots\cup B_{j-1}$.
\end{enumerate}
The empty sequence satisfies these conditions.  Put
$B=\bigcup_j B_j$, and let $U$ be the union of the candidate palettes
and all palettes in $B$.  Each element of $B$ intersects exactly its
associated candidate; these intersections are disjoint because matching
palettes are disjoint.  Since $t\le|B|$,
\[
 |U|\le rt+r|B|-|B|\le(2r-1)|B|.
\]
Let $J$ contain the root and all families represented by $B$, so that
$|J|=|B|+1$.  The hypothesis guarantees a palette $a$ from a family in
$J$ that avoids $U$.

If $a$ intersects an element of $M$, append $a$ and all its blockers
to the sequence.  Its blockers are new, because $a$ avoids $U$.
A new blocker cannot intersect an earlier candidate, since all blockers
of those candidates were already in $B$.  Thus every invariant persists.

Otherwise $a$ is unblocked.  If its family is the root, insert $a$ into
$M$ and finish.  If not, its family has a unique matching element
$b\in B_j$.  Replace $b$ in $M$ by $a$, discard all sequence entries
after $j$, and remove $b$ from $B_j$.  This is a valid matching swap,
and earlier candidates avoid the new matching element $a$.
If $B_j$ is still nonempty, resume the procedure.  If it is empty,
remove entry $j$ and treat $a_j$ as the next unblocked palette.
By the third invariant, its family is the root or has its matching
element in an earlier blocker set.  Repeat the same swap rule.
This cascade moves strictly backwards, so it either finishes at the
root or stops with an admissible sequence having only nonempty blocker
sets.  The earlier invariants remain valid throughout: a palette
inserted into $M$ avoids all earlier candidates, and no earlier candidate
depends on the removed blocker; any later candidate that did depend
on it was discarded.

To prove termination, associate with every admissible sequence the
length-$n$ integer vector
\[
 (|B_1|,\ldots,|B_t|,n,\ldots,n).
\]
There are at most $n-1$ matching elements, and $t\le n-1$.
Appending a blocker set strictly decreases this vector in lexicographic
order, because the new entry is at most $n-1$.
A swap and its backwards cascade, unless they finish, leave a prefix
unchanged and decrease the blocker count at the next retained entry;
truncation affects only later coordinates.  They therefore strictly
decrease the vector as well.  Only finitely many such vectors exist.
The procedure must finish the augmentation, proving the induction step.
\end{proof}

\section{An explicit coefficient of \texorpdfstring{$16.5$}{16.5}}
\label{app:constant}

Write $d=\Delta$ throughout this appendix. The elementary proof of Theorem~\ref{thm:linearordinary} gives $512d$ colors.
The following refinement preserves its asymptotic round complexity.  Its
purpose is to record a concrete coefficient; no optimality is claimed.

We prove six constant-parameter palette reductions, as used in Theorem~\ref{thm:smallconstant}.  In every row of
Table~\ref{tab:constant-parameters}, an input palette of size
$\lfloor Ld\rfloor$ is replaced by one of size $\lfloor Rd\rfloor$.
Every sampled candidate has exactly six elements.

\begin{table}[ht]
\centering
\begin{tabular}{c r r r}
\toprule
Row & $L$ & $R$ & Candidates per family $m$\\
\midrule
1 & $512$ & $100$ & $500$\\
2 & $100$ & $20$ & $3500$\\
3 & $20$ & $18$ & $6200$\\
4 & $18$ & $17$ & $10400$\\
5 & $17$ & $83/5$ & $10200$\\
6 & $83/5$ & $33/2$ & $10500$\\
\bottomrule
\end{tabular}
\caption{Six reductions with a common existence proof.}
\label{tab:constant-parameters}
\end{table}

Fix a row, and suppose $d\ge6000$.  Put
\[
 M=\lfloor Ld\rfloor,\qquad Q=\lfloor Rd\rfloor,\qquad
 S=R-\frac1{1000},\qquad \kappa=11.
\]
Independently sample $m$ uniform six-element subsets of $[Q]$ for each of
$M$ families.  Repetitions between samples are allowed.  The choice of
$S$ handles both rounding and sampling without replacement:
\begin{equation}\label{eq:constant-rounding}
 Q-5\ge Rd-6\ge Sd.
\end{equation}

There are two kinds of bad events.  A failure of robust Hall yields $k$
distinct families and a $(k-1)$-element set containing a candidate from
each family, for some $1\le k\le d$.  The probability of such a witness,
already summed over the choices of families and of the containing set,
is at most
\begin{equation}\label{eq:constant-hall-count}
 A_k\le \binom{M}{k}\binom{Q}{k-1}
       \left[m\left(\frac{k}{Sd}\right)^6\right]^k.
\end{equation}
For $k=1$ the event is impossible, but including its upper bound is harmless.
For disjoint palette choices, we impose the stronger sufficient condition
that no $k$ families have all their candidates met by an $11k$-element
set, for any $1\le k\le d$.  This implies the hypothesis of
Theorem~\ref{lem:haxell}: a smaller forbidden transversal can be enlarged
because $Q\ge11d$.  Write
\[
 h(x)=1-\left(1-\frac{11x}{S}\right)^6,\qquad 0<x\le1.
\]
For a fixed set of size $11k$, the probability that a uniform six-element
subset meets it is
\[
 1-\prod_{j=0}^5\left(1-\frac{11k}{Q-j}\right)
 \le h(k/d),
\]
by \eqref{eq:constant-rounding}.  Thus the second bad-event probability is
\begin{equation}\label{eq:constant-cover-count}
 B_k\le\binom{M}{k}\binom{Q}{11k}h(k/d)^{mk}.
\end{equation}

Here is a uniform bound on both sums.  Let
$H(z)=-z\log z-(1-z)\log(1-z)$ denote binary entropy, with natural
logarithms, and define
\begin{align*}
 F_H(x)&=R H(x/R)+L H(x/L)+x\log m+6x\log(x/S),\\
 F_C(x)&=L H(x/L)+R H(11x/R)+mx\log h(x),\\
 p&=h(1),\qquad
 \alpha=\frac{h'(1)}{h(1)},\qquad
 \beta=m\alpha-12,\\
 a&=\frac{e^2 LRm}{S^6},\qquad
 b=eL\left(\frac{eR}{11}\right)^{11}p^m.
\end{align*}
The following numerical bounds hold for all six rows:
\begin{equation}\label{eq:constant-margins}
 F_H(1)<-\frac9{1000},\quad F_C(1)<-\frac3{25},\quad
 a<2,\quad \beta>160,\quad \log b<6.
\end{equation}
A rational certificate for these bounds is given below.

For completeness, the entropy estimate used here follows directly from
the binomial theorem: for $0<t<N$, the $t$th summand in
$(t/N+(1-t/N))^N=1$ gives
$\binom Nt\le\exp(NH(t/N))$; the endpoints follow by continuity.
For fixed $t$, the function $uH(t/u)$ increases with $u$, since its
 derivative is $-\log(1-t/u)>0$.  Hence the same upper bound is valid with
any real upper bound on the population.  Also
$\binom Q{k-1}\le\binom Qk$ for $k\le d$, since $Q\ge2d$.
Applying these facts to \eqref{eq:constant-hall-count} and
\eqref{eq:constant-cover-count} gives
\begin{equation}\label{eq:constant-entropy-bounds}
 A_k\le e^{dF_H(k/d)},\qquad B_k\le e^{dF_C(k/d)}.
\end{equation}

Both ratios $F_H(x)/x$ and $F_C(x)/x$ increase on $(0,1]$.
Indeed, differentiation gives
\begin{align*}
 x^2\left(\frac{F_H(x)}x\right)'
 &=R\log(1-x/R)+L\log(1-x/L)+6x,\\
 x^2\left(\frac{F_C(x)}x\right)'
 &=L\log(1-x/L)+R\log(1-11x/R)
   +mx\frac{xh'(x)}{h(x)}.
\end{align*}
Writing $y=1-11x/S$ shows that the elasticity
\[
 \frac{xh'(x)}{h(x)}
 =\frac{6}{1+y^{-1}+\cdots+y^{-5}}
\]
is decreasing, and hence is at least $\alpha$.
The inequality $-\log(1-u)\le u/(1-u)$, together with
$L,R\ge33/2$, bounds the two negative terms in the first derivative
below by $-2x$ each.  In the second derivative they are bounded below
by $-2x$ and $-33x$, respectively.  Since
$m\alpha=\beta+12>172$, both derivatives are positive.
It follows from \eqref{eq:constant-margins} and
\eqref{eq:constant-entropy-bounds} that for $k>d/2$,
\[
 A_k\le e^{-9k/1000},\qquad B_k\le e^{-9k/1000}.
\]
Each large-witness sum is therefore smaller than
\[
 \frac{e^{-27}}{1-e^{-9/1000}}
 \le\frac{1009}{9\cdot2^{27}}<\frac1{100}.
\]
Here $d\ge6000$, $e>2$, and
$1-e^{-u}\ge u/(1+u)$ were used.

Small witnesses require the factors depending on $k/d$.  The bound
$\binom Nt\le(eN/t)^t$ applied to the same counting formulas gives
\[
 A_k\le\left[a(k/d)^4\right]^k.
\]
Integrating the elasticity inequality gives $h(x)\le px^\alpha$, so
\[
 B_k\le\left[b(k/d)^\beta\right]^k.
\]
For $k\le d/2$, the first bracket is less than $2/16=1/8$.
The second is less than $e^6 2^{-160}<1/8$.
Each small-witness sum is at most $\sum_{k\ge1}8^{-k}=1/7$.
Consequently the probability of any obstruction is less than
\[
 \frac27+\frac2{100}<1.
\]
A suitable table exists in every row.  Theorem~\ref{lem:haxell} and
Lemma~\ref{lem:protocol} turn it into the stated palette reduction.

\paragraph{Self-contained numerical certificate.}
All entries of Table~\ref{tab:constant-certificate} are rational bounds;
a terminating recipe to check them follows.  In its last two columns,
$A=9LRm/S^6$ is a rational upper bound on $a$ because $e<3$, and
$\beta=m\alpha-12$ is itself rational.  The entries in the two middle
columns are open intervals containing the indicated endpoint values.
\begin{table}[ht]
\centering
\small
\begin{tabular}{c c c r r}
\toprule
Row & $F_H(1)$ & $F_C(1)$ & $A<$ & $\beta>$\\
\midrule
1 & $(-8.578853,-8.578852)$ & $(-301.671413,-301.671412)$ & $0.001$ & $354.334$\\
2 & $(-0.243118,-0.243117)$ & $(-9.810845,-9.810844)$ & $0.985$ & $202.859$\\
3 & $(-0.777223,-0.777222)$ & $(-5.473060,-5.473059)$ & $0.591$ & $190.828$\\
4 & $(-0.084094,-0.084093)$ & $(-5.209509,-5.209508)$ & $1.188$ & $209.447$\\
5 & $(-0.044046,-0.044045)$ & $(-0.619635,-0.619634)$ & $1.239$ & $165.355$\\
6 & $(-0.009584,-0.009583)$ & $(-0.121551,-0.121550)$ & $1.284$ & $160.982$\\
\bottomrule
\end{tabular}
\caption{Outward rational bounds for the six parameter rows.}
\label{tab:constant-certificate}
\end{table}

For a positive rational number $u$, choose an integer $j$ so that
$1\le v=2^{-j}u<2$, and put $z=(v-1)/(v+1)$.  With $N=40$, define
\[
 P(v)=2\sum_{i=0}^{N-1}\frac{z^{2i+1}}{2i+1},\qquad
 E(v)=\frac{2z^{2N+1}}{(2N+1)(1-z^2)}.
\]
The geometric expansion of $1/(1-z^2)$, integrated term by term on
$[0,z]$, proves
\[
 P(v)\le\log v\le P(v)+E(v).
\]
Apply the same formula at $v=2$ to enclose $\log2$, and use
$\log u=j\log2+\log v$.  When multiplying an interval by a negative
number, interchange its endpoints.  Evaluate $H$, $F_H(1)$, and $F_C(1)$
from their displayed formulas using interval addition and rational scalar
multiplication.  This produces the open enclosures in
Table~\ref{tab:constant-certificate}; every operation is rational, and
the error formula specifies all discarded terms.  The inequalities for
$A$ and $\beta$ are direct rational substitutions.

The remaining estimate $\log b<6$ does not need another numerical check.
Expanding the entropy terms gives
\begin{align*}
 \log b-F_C(1)
 &=1+(L-1)\log(1-1/L)\\
 &\quad+11+(R-11)\log(1-11/R).
\end{align*}
The first line is less than $1/L\le2/33$.
The function $11+(R-11)\log(1-11/R)$ decreases with $R$, since its
derivative is $\log(1-11/R)+11/R<0$.  At $R=33/2$ it equals
$11-(11/2)\log3<11/2$, because $e<3$.
Since $F_C(1)<0$, these estimates imply $\log b<6$ and complete
\eqref{eq:constant-margins}.


\section{Elementary vertex-coloring tools}\label{app:defective}
\begin{lemma}[Defective color reduction]\label{lem:defective-elementary}
Let a graph of maximum degree at most $D\ge1$ have a proper vertex
$M$-coloring.  For every real $p\ge1$, a deterministic algorithm computes
a vertex coloring with defect at most $D/p$ and
$O(p^2\log(2p))$ colors in $O(1+\log^* M)$ rounds.
The constants in both bounds are absolute.
\end{lemma}

\begin{proof}
We give a slightly less economical version of the color-reduction method
of Kuhn~\cite[Section~4]{Kuhn}; the extra logarithm in the palette will be
irrelevant to our applications.  A coloring has defect $b$ if each vertex
has at most $b$ neighbors with its color.

We first prove a one-round reduction.  Suppose the current coloring has
$m\ge2$ colors and defect at most $b$, and choose an integer $q\ge1$.
Put $a=\lceil8q\ln m\rceil$.  There exist functions
$f_i:[a]\to[4q]$, one for each old color $i\in[m]$, such that distinct
functions agree at at most $a/q$ arguments.  Indeed, choose all their
values independently and uniformly.  For two fixed indices, the
probability of agreement at more than $a/q$ arguments, with
$t=\lfloor a/q\rfloor+1$, is zero if $t>a$, and otherwise at most
\[
 \binom at(4q)^{-t}
 \le \left(\frac{3a}{4qt}\right)^t
 \le (3/4)^t
 \le (3/4)^{8\ln m}.
\]
A union bound over fewer than $m^2$ pairs gives probability less than one,
since $8\ln(4/3)>2$.  Fix such a table by finite exhaustive search.

After receiving its neighbors' old colors, a vertex of old color $i$
chooses an argument $x$ minimizing the number of neighbors of color $j$
for which $f_j(x)=f_i(x)$, and outputs $(x,f_i(x))$.
Neighbors with old color $i$ contribute at most $b$ to this count at
every argument.  Averaged over the arguments, all other neighbors
contribute at most $D/q$.  Consequently the new defect is at most
$b+D/q$.  A neighbor with the same new color is necessarily included
in the count at the chosen argument, even though it chooses its own
argument independently.  Writing $L=\log_2 m\ge1$, the new number of
colors is at most
\begin{equation}\label{eq:defective-one-step}
 4qa\le64q^2L.
\end{equation}

Let $K=2^{15}$.  Starting from the proper coloring, repeat this reduction
while $m>(Kp)^{32}$, each time taking
$q=\lceil16p(\log_2m)^2\rceil$.  The new palette bound $m'$ satisfies
\[
 m'\le Kp^2L^5,\qquad
 L':=\log_2m'\le A+5\log_2L,
 \qquad A=15+2\log_2p.
\]
Whenever a reduction is performed, $L>32(15+\log_2p)\ge480$.
Thus $A<L/16$ and $\log_2L\le L/16$, whence $L'<L/2$.
If the successive logarithms before these reductions are
$L_0,\ldots,L_{t-1}$, their reciprocal squares have sum at most $4/3$:
read the sequence backwards and use $L_i>2L_{i+1}$ and
$L_{t-1}\ge1$.  The total increase in defect is therefore at most
\[
 \sum_{i<t}\frac{D}{16pL_i^2}\le\frac{D}{12p}.
\]

There are $O(1+\log^*M)$ reductions, uniformly in $p$.
To see this explicitly, if $5\log_2L\le A$, then
$L'\le2A<32(15+\log_2p)$ and the loop stops.
Every other step obeys $L'\le10\log_2L$.
Above a fixed absolute threshold, two applications of
$x\mapsto10\log_2x$ are bounded by one application of $x\mapsto\log_2x$.
For example, for $x\ge2^{1024}$ this follows from
$10\log_2(10z)\le z$ for $z=\log_2x\ge1024$.
Below that threshold, the preceding halving bound allows only an
absolute constant number of further iterations.

At the end $m\le(Kp)^{32}$.  Apply the one-round reduction once more
with $q=\lceil2p\rceil$.  Its defect contribution is at most $D/(2p)$,
so the final defect is less than $D/p$.  By
\eqref{eq:defective-one-step}, its palette has size
\[
 O(p^2\log m)=O(p^2\log(2p)).
\]
If the initial proper coloring has at most one color, the graph is
edgeless and requires no reduction.  This completes the proof.
\end{proof}

\begin{corollary}[Initial proper vertex coloring]\label{cor:initial-vertex-coloring}
A graph of maximum degree at most $D\ge1$, with a proper $M$-vertex coloring, admits a deterministic proper $O(D^2)$-vertex coloring in $O(1+\log^* M)$ rounds. In particular, polynomial-size unique identifiers on an $n$-vertex graph give such a coloring in $O(1+\log^*n)$ rounds.
\end{corollary}

\begin{proof}
This is Linial's color-reduction theorem~\cite{Linial}; the following
derivation supplies the version needed here.  Starting from the proper $M$-coloring, apply
Lemma~\ref{lem:defective-elementary} with $p=2D$.
The resulting defect is less than $1/2$, hence zero, and its palette
has size at most $C D^2\log(2D)$ for an absolute constant $C$.
Choose a sufficiently large fixed constant $c$, and a prime
$cD<q<2cD$, as guaranteed by Bertrand's postulate.
Taking $c$ large enough ensures both $q>2D$ and
$q^3\ge CD^2\log(2D)$ for every integer $D\ge1$.
Encode the old colors as distinct polynomials of degree at most two
over the field $\mathbb F_q$.
After learning its neighbors' polynomials, each vertex chooses a field
element $x$ at which its polynomial differs from every neighbor's
polynomial.  Such an $x$ exists, since each neighbor forbids at most
two arguments and $q>2D$.  Output $(x,P(x))$.
If adjacent vertices choose the same argument, their values differ;
otherwise their first coordinates differ.  This is therefore a proper
$q^2=O(D^2)$-coloring, obtained in one additional round. For the identifier corollary take $M=n^{O(1)}$.
\end{proof}

\section{Constant-time coloring from large lists}\label{app:fhk}
\begin{lemma}[Constant-round list coloring with large lists]
\label{lem:fhk-elementary}
There is an absolute constant $B$ such that a graph of maximum degree
$D\ge2$, supplied with a proper $s$-vertex-coloring, admits a
deterministic list-coloring algorithm in constant time whenever every
list has at least
\[
 \max\{s,\lceil B D^2\ln(2D)\rceil\}
\]
colors.  The color labels may be arbitrary finite strings.  For
example, $B=2^{20}$ suffices.
\end{lemma}

\begin{proof}
We prove the needed special case of the conflict-coloring method of
Fraigniaud, Heinrich, and Kosowski~\cite{FHK}, with deliberately generous
constants.  A conflict-coloring instance has the common local list
$[l]$ at each vertex.  Each edge specifies forbidden pairs of local
colors; for each color at either endpoint, at most $b$ colors at the
other endpoint are forbidden.  Write $R=l/b$.

We first describe one simplification step, valid when $R\ge96D^2$.
Set
\[
 k=\left\lfloor\frac{l}{16bD}\right\rfloor,
 \qquad t=\left\lfloor\frac{k}{D}\right\rfloor-1,
 \qquad N=\binom lk,
 \qquad Q=\binom{kb}{t}\binom l{k-t}.
\]
Candidates for the new local colors are the $k$-subsets of $[l]$.
Declare candidates $S$ and $T$ at adjacent vertices incompatible if
more than $t$ elements of $S$ conflict with at least one element of $T$,
or more than $t$ elements of $T$ conflict with at least one element of
$S$.  For a fixed $S$, its elements conflict with at most $kb$ colors
at the other endpoint.  Thus there are at most $Q$ candidates $T$
containing at least $t$ such colors: choose $t$ of these colors and
then the remaining $k-t$ elements.  Counting the other condition with
the endpoints exchanged shows that the bipartite incompatibility graph
between all $N$ candidates at either endpoint has at most $2NQ$ edges.

At a vertex discard every candidate with more than $4DQ$
incompatibilities along any one incident edge.  For each edge at most
$N/(2D)$ candidates are discarded, so at least $N/2$ remain.  Retain
the first $l'=\lfloor N/2\rfloor$ in lexicographic order and renumber
them by $[l']$.  Exchange these retained lists with neighbors.
The new conflict instance has conflict degree at most $b'=4DQ$;
all of its local input, including the renumbered forbidden pairs,
is obtained in one communication round from the previous instance.

Conversely, any solution to the new instance specifies a candidate
$S_v$ at each vertex.  After one exchange of these candidates, vertex
$v$ chooses an element of $S_v$ conflicting with no element of any
neighbor's candidate.  Each neighbor forbids at most $t$ elements
of $S_v$, and $Dt<k$, so such an element exists.  The resulting
colors solve the preceding conflict instance.

We record how much the list-to-conflict ratio improves.  Our assumptions
give $k\ge4D$, $t\ge k/(2D)$, and $l-k\ge15bDk$.
Using $\binom{kb}{t}\le(ekb/t)^t$ and $l'\ge N/3$, we obtain
\[
 \frac{l'}{b'}
 \ge\frac1{12D}\frac{\binom lk}
          {\binom{kb}{t}\binom l{k-t}}
 \ge\frac1{12D}
       \left(\frac{(l-k)t}{ek^2b}\right)^t
 \ge\frac{2^t}{16D}.
\]
Indeed, the base in parentheses is at least $15/(2e)>2$.
Also
$t\ge R/(16D^2)-3\ge R/(32D^2)$.
Consequently each simplification guarantees
\begin{equation}\label{eq:conflict-amplify}
 R'\ge\frac1{16D}\exp\left(\frac{R}{64D^2}\right).
\end{equation}

Now truncate the original lists to the common length
$l=\max\{s,\lceil2^{20}D^2\ln(2D)\rceil\}$.
In one round exchange the ordered original lists and local port numbers.
At each vertex
replace its labels by their local ranks in $[l]$ and record, for each
neighbor, the pairs of ranks representing equal actual colors.
This is a conflict instance with $b=1$.  The original labels will only
be used again to decode the eventual answer.

Apply the simplification twice, obtaining ratios $R_1$ and $R_2$.
Equation~\eqref{eq:conflict-amplify} gives
\[
 R_1\ge l^8,
 \qquad R_2\ge\exp(l^6).
\]
Here are elementary estimates covering all allowed parameters.
The function $x/(64D^2)-8\ln x$ is increasing for
$x\ge2^{20}D^2\ln(2D)$.  At its left endpoint,
\[
 8\ln x+\ln(16D)
 \le108\ln(2D)
 <2^{14}\ln(2D)=x/(64D^2),
\]
where we used
$\ln(2^{20})\le10\ln(2D)$,
$\ln D\le\ln(2D)$,
$\ln\ln(2D)\le\ln(2D)$, and
$\ln(16D)\le4\ln(2D)$.
This proves the first bound.  Since $D^2\le l$ and $l\ge128$,
$l^8/(64D^2)-\ln(16D)\ge l^6$, proving the second.
In particular the hypotheses of both simplifications hold.

It remains to solve the twice-simplified instance without further
communication.  There are only finitely many possible normalized
local inputs at this point.  We give a bound that does not depend on
the original color names.  An initial normalized local input consists
of its degree, its incident ports and opposite port numbers, and at
most $D$ Boolean $l\times l$ conflict matrices.  Use $D$ port slots, each storing a presence bit, an opposite-port number in $\lceil\log_2(D+1)\rceil$ bits, and an $l\times l$ conflict matrix. Since $l\ge128$ and $D^2\le l$, this uses at most $2Dl^2$ bits.  After two simplification rounds a local
input is a deterministic function of a depth-two unfolded port-labelled
neighborhood.  Such a neighborhood has at most $2D^2$ vertices.
Including its port incidences and the normalized initial inputs,
it has a fixed-length encoding of $8D^3l^2$ bits, padding unused fields.
Hence the set
$\mathcal I$ of possible twice-simplified local inputs satisfies
\[
 |\mathcal I|\le2^{8D^3l^2}.
\]
All vertices can enumerate a common finite superset of these inputs
by enumerating such neighborhood encodings, performing the two
local transformations, and discarding malformed encodings, encodings with initial conflict degree greater than one, and encodings for which a required retained list is too short. Retain the resulting local inputs whose conflict matrices obey the prescribed conflict-degree bound. Every actual input survives these checks.

Pair each type in $\mathcal I$ with a proper vertex color in $[s]$.
In a fixed order, assign an output from the common twice-simplified
list to every such pair.  When assigning a pair, avoid any color
that conflicts, along any local port, with the output already assigned
to any earlier pair.  Each earlier pair excludes at most $Db_2$
colors, where $b_2$ is the new conflict-degree bound.  This greedy
table exists because
\[
 D s|\mathcal I|<\exp(l^6)\le R_2=l_2/b_2.
\]
For the strict inequality, use $s\le l$, $D^2\le l$, and $l\ge128$:
the logarithm of its left side is at most
$2\ln l+8D^3l^2\ln2\le9l^4<l^6$.
Adjacent vertices have different proper vertex colors, hence different
type--color pairs.  The later pair in the fixed order avoids the
earlier pair's output on their actual edge.  Looking up one's pair
therefore gives a valid conflict-coloring without communication.

Finally invert the two simplifications, one round each, and translate
local ranks back to actual input colors.  This uses a constant number
of communication rounds and always returns a color in the original
list.  The construction is deterministic: all tables are finite and
chosen canonically, with unrestricted local computation.
\end{proof}

\phantomsection
\addcontentsline{toc}{section}{References}
\begingroup\small
\begin{thebibliography}{99}

\bibitem{Hall}
Philip Hall.
On representatives of subsets.
\emph{Journal of the London Mathematical Society}, 10:26--30, 1935.
\href{https://doi.org/10.1112/jlms/s1-10.37.26}{doi:10.1112/jlms/s1-10.37.26}.

\bibitem{Haxell}
Penny E. Haxell.
A condition for matchability in hypergraphs.
\emph{Graphs and Combinatorics}, 11(3):245--248, 1995.
\href{https://doi.org/10.1007/BF01793010}{doi:10.1007/BF01793010}.

\bibitem{Annamalai}
Chidambaram Annamalai.
Finding perfect matchings in bipartite hypergraphs.
In \emph{Proceedings of the ACM--SIAM Symposium on Discrete Algorithms
(SODA)}, 2016.
Full version: \href{https://arxiv.org/abs/1509.07007v2}{arXiv:1509.07007v2}, 2016.

\bibitem{BGIKS}
Radu Berinde, Anna C. Gilbert, Piotr Indyk, Howard Karloff, and Martin J. Strauss.
Combining geometry and combinatorics: A unified approach to sparse signal recovery.
\href{https://arxiv.org/abs/0804.4666}{arXiv:0804.4666}, 2008.

\bibitem{Linial}
Nathan Linial.
Locality in distributed graph algorithms.
\emph{SIAM Journal on Computing}, 21(1):193--201, 1992.
\href{https://doi.org/10.1137/0221015}{doi:10.1137/0221015}.

\bibitem{Kuhn}
Fabian Kuhn.
Weak graph colorings: Distributed algorithms and applications.
In \emph{Proceedings of the ACM Symposium on Parallelism in Algorithms
and Architectures (SPAA)}, pages 138--144, 2009.
\href{https://people.csail.mit.edu/fkuhn/publications/spaa09b.pdf}{Author's version}.

\bibitem{BEK}
Leonid Barenboim, Michael Elkin, and Fabian Kuhn.
Distributed $(\Delta+1)$-coloring in linear (in $\Delta$) time.
\emph{SIAM Journal on Computing}, 43(1):72--95, 2014.
\href{https://doi.org/10.1137/12088848X}{doi:10.1137/12088848X}.

\bibitem{BKO}
Alkida Balliu, Fabian Kuhn, and Dennis Olivetti.
Distributed edge coloring in time quasi-polylogarithmic in Delta.
In \emph{Proceedings of the ACM Symposium on Principles of Distributed
Computing (PODC)}, pages 289--298, 2020.
Full version: \href{https://arxiv.org/abs/2002.10780}{arXiv:2002.10780}.

\bibitem{BBKO}
Alkida Balliu, Sebastian Brandt, Fabian Kuhn, and Dennis Olivetti.
Distributed edge coloring in time polylogarithmic in $\Delta$.
In \emph{Proceedings of the ACM Symposium on Principles of Distributed
Computing (PODC)}, 2022.
Full version: \href{https://arxiv.org/abs/2206.00976}{arXiv:2206.00976}.
Journal version: \emph{SIAM Journal on Computing}, 55(4), 2026,
\href{https://doi.org/10.1137/24M1715866}{doi:10.1137/24M1715866}.

\bibitem{FHK}
Pierre Fraigniaud, Marc Heinrich, and Adrian Kosowski.
Local conflict coloring.
In \emph{Proceedings of the IEEE Symposium on Foundations of Computer
Science (FOCS)}, 2016.
Full version: \href{https://arxiv.org/abs/1511.01287v2}{arXiv:1511.01287v2}, 2017.

\bibitem{MT}
Yannic Maus and Tigran Tonoyan.
Local conflict coloring revisited: Linial for lists.
In \emph{34th International Symposium on Distributed Computing (DISC 2020)},
volume 179 of \emph{LIPIcs}, pages 16:1--16:18, 2020.
\href{https://doi.org/10.4230/LIPIcs.DISC.2020.16}{doi:10.4230/LIPIcs.DISC.2020.16}.

\bibitem{CHLPU}
Yi-Jun Chang, Qizheng He, Wenzheng Li, Seth Pettie, and Jara Uitto.
Distributed edge coloring with small palettes and a special case of the
constructive Lov\'asz local lemma.
\href{https://arxiv.org/abs/1708.04290v3}{arXiv:1708.04290v3}, 2026.
Preliminary version: \emph{The complexity of distributed edge coloring with
small palettes}, in \emph{Proceedings of SODA}, 2018.

\bibitem{Brandt}
Sebastian Brandt.
An automatic speedup theorem for distributed problems.
In \emph{Proceedings of the ACM Symposium on Principles of Distributed
Computing (PODC)}, 2019.
Full version: \href{https://arxiv.org/abs/1902.09958}{arXiv:1902.09958}.

\end{thebibliography}
\endgroup

\end{document}
